% ═══════════════════════════════════════════════════════════════════════════════
% PLANTEAMIENTO DEL PROBLEMA - INSTITUTO ATENEO DE COLIMA
% Estudio Documental: Corrupción en México (2006-2012)
% ═══════════════════════════════════════════════════════════════════════════════

\nota{
  Este capítulo plantea el problema de investigación desde una perspectiva 
  documental basada en evidencia oficial (ASF, SCJN, CNDH, CIDH, DOJ) y 
  doctrina académica (Casar y Cejudo 2024; Merino 2024). Se enfatiza la 
  brecha entre normatividad anticorrupción y aplicación efectiva.
}

\section{Contexto institucional y normativo}

A pesar de contar con una extensa arquitectura constitucional, administrativa y penal ---integrada por el artículo 113 constitucional, el Sistema Nacional Anticorrupción (SNA), la Ley General de Responsabilidades Administrativas (LGRA), el Título Décimo del Código Penal Federal y convenciones vinculantes ante la ONU, la OEA y la OCDE---, México se mantuvo en niveles críticos de corrupción percibida y experimentada \citep{cpeum,lgra,casar2024,merino2024}. En el IPC 2025, difundido en febrero de 2026 por Transparencia Internacional, el país obtuvo 27 de 100 puntos y ocupó el lugar 141 de 182 naciones evaluadas, último puesto entre los países miembros de la OCDE y siete puntos por debajo de la calificación de 34 puntos registrada al concluir 2012 \citep{transparencia2026,mcci2026b}.

Durante el sexenio de Felipe Calderón Hinojosa (2006--2012), la fragilidad de los contrapesos horizontales y verticales favoreció el desarrollo de esquemas de corrupción de alto calado cuyos efectos rebasaron los límites del periodo presidencial \citep{astorga2015,lopezayllon2009}. En primer lugar, la concentración de mando, presupuesto e inteligencia en la Secretaría de Seguridad Pública (SSP) permitió la colusión de su titular, Genaro García Luna, con el Cártel de Sinaloa y la operación de una red de desvío de recursos mediante contratos tecnológicos y empresas offshore \citep{doj2024garcialuna,uif2025,uif2026}. En segundo término, proyectos estratégicos impulsados sin candados técnicos vinculantes, como la Refinería Bicentenario en Tula, comprometieron endeudamiento subnacional y miles de millones de pesos federales sin concluir una infraestructura productiva \citep{asf2015,asf2016}. En tercer lugar, la expansión acelerada de guarderías subrogadas del IMSS bajo estándares flexibilizados, ausencia de inspección eficaz y conflictos de interés culminó en el incendio de la Guardería ABC, con un saldo de 49 niñas y niños fallecidos \citep{cndh2009,scjn2010}.

A estos problemas estructurales se sumó una distorsión persistente en el espacio público: los encuadres mediáticos (framing) y la polarización partidista suelen simplificar fenómenos complejos en símbolos virales, mezclando delitos acreditados, negligencia regulatoria, ineficiencia administrativa y afirmaciones sin sustento sanitario o financiero \citep{solamorales2019}. Ante este escenario, resultó necesario examinar de manera sistemática y verificable cómo se configuraron los mecanismos de corrupción en los tres casos del periodo 2006--2012 y cuál fue su impacto real sobre las instituciones, el erario y los derechos humanos.

\section{Evaluación de las deficiencias en el conocimiento del problema}

Conforme a los criterios de planteamiento establecidos por Hernández Sampieri et al. (2014), la revisión inicial de la literatura evidenció una brecha relevante en el conocimiento acumulado sobre la corrupción del sexenio 2006--2012. Por un lado, existen expedientes jurídicos y resoluciones judiciales especializadas (como las sentencias penales y civiles estadounidenses contra García Luna o los fallos de la SCJN y la CIDH sobre Guardería ABC); por otro, auditorías técnicas y financieras de la ASF sobre obra pública y pasivos estatales; y, en una vertiente separada, estudios de comunicación política acerca del encuadre mediático (framing) de los escándalos públicos. Sin embargo, se identificó una escasez de estudios integradores que articulen simultáneamente esas tres dimensiones ---evidencia jurídico-judicial, impacto financiero auditado y análisis de encuadre mediático--- incorporando las resoluciones emitidas entre 2024 y noviembre de 2026. Esta investigación se propuso cubrir dicha deficiencia mediante un análisis comparado de casos múltiples.

\section{Consecuencias éticas de la investigación}

De acuerdo con Hernández Sampieri et al. (2014), todo planteamiento científico exige evaluar las consecuencias éticas derivadas de su ejecución y difusión. Dado que el estudio abordó delitos graves, procesos judiciales activos y una tragedia que involucró a menores de edad (Guardería ABC), se aplicaron tres principios éticos estrictos: primero, el uso exclusivo de fuentes públicas, oficiales y verificables (resoluciones judiciales, auditorías, recomendaciones de derechos humanos y literatura académica), evitando especulaciones o señalamientos carentes de soporte documental; segundo, el respeto irrestricto a la dignidad de las víctimas y sus familias, excluyendo descripciones sensacionalistas o cualquier forma de revictimización; y tercero, el mantenimiento de un balance probatorio objetivo que registró explícitamente las posturas de descargo de los actores públicos señalados y la continuidad transexenal de diversos entramados, evitando sesgos partidistas.